\section{Кратность корней}

Многочлены $f,g\in F[X]$, не имеющие общего множителя в $F[X]$, могли бы, казалось бы, приобрести общий множитель в
$\Om[X]$ для некоторого расширения $\Om\supset F$. На самом деле этого не происходит: наибольший общий делитель не
меняется при расширении поля. Здесь и далее $\gcd(f,g)$~--- \emph{унитарный} наибольший общий делитель многочленов $f,g$,
не равных нулю одновременно (\rf{1.8}; при $g=0$ полагаем $\gcd(f,0)=\gcd(0,f)=c^{-1}f$, где $c$~--- старший
коэффициент $f$, а если $f$ и $g$~--- ненулевые константы, то $\gcd(f,g)=1$).

\begin{proposition}[2.17]
Пусть $f,g\in F[X]$ и $\Om$~--- расширение $F$. Если $r(X)$~--- наибольший общий делитель $f$ и $g$, вычисленный в
$F[X]$, то $r$~--- также наибольший общий делитель $f$ и $g$ в $\Om[X]$. В частности, различные унитарные неприводимые
многочлены из $F[X]$ не приобретают общего корня ни в каком расширении $F$.
\end{proposition}

\begin{proof}
Пусть $r_F$ и $r_\Om$~--- унитарные наибольшие общие делители $f$ и $g$ в $F[X]$ и в $\Om[X]$ соответственно.
Многочлен $r_F$ делит $f$ и $g$ в $F[X]$, а значит, и в $\Om[X]$; поэтому $r_F\mid r_\Om$ в $\Om[X]$. С другой
стороны, по алгоритму Евклида (\rf{1.8}) существуют $a,b\in F[X]$ (в вырожденных случаях это очевидно: при $g=0$ годятся
$a=c^{-1}$, $b=0$, где $c$~--- старший коэффициент $f$, при $f=0$~--- симметрично, а при $f,g\in F^\times$~--- $a=f^{-1}$,
$b=0$), для которых
\[
a(X)f(X)+b(X)g(X)=r_F(X),
\]
и это равенство выполняется и в $\Om[X]$. Многочлен $r_\Om$ делит левую часть, значит, $r_\Om\mid r_F$ в $\Om[X]$.
Два унитарных многочлена, делящих друг друга, равны: $r_F=r_\Om$.

Пусть теперь $f\ne g$~--- унитарные неприводимые многочлены из $F[X]$. Делители $f$ в $F[X]$~--- это константы и
многочлены $cf$, $c\in F^\times$; если бы $f$ делил $g$, то из неприводимости $g$ следовало бы $g=cf$, т.\,е.\ $g=f$ (оба
унитарны). Значит, $\gcd(f,g)=1$ в $F[X]$, а тогда и в $\Om[X]$. Если бы $f$ и $g$ имели общий корень $\alpha\in\Om$, то
$X-\alpha$ был бы их общим делителем в $\Om[X]$ (\rf{1.7}), что невозможно.
\end{proof}

Таким образом, можно говорить о наибольшем общем делителе $f$ и $g$, не указывая поля. Грубо говоря, \emph{различные
неприводимые многочлены имеют различные корни}, в каком бы расширении поля $F$ мы ни искали эти корни. Так же от поля не
зависит и делимость.

\begin{lemma}[2.17a]
Пусть $f,g\in F[X]$, $g\ne0$, и $\Om\supset F$. Если $g\mid f$ в $\Om[X]$, то $g\mid f$ в $F[X]$.
\end{lemma}

\begin{proof}
Разделим с остатком в $F[X]$: $f=gq+r$, где $r=0$ или $\deg r<\deg g$ (\rf{1.6}). По условию $f=gq_1$ с $q_1\in\Om[X]$.
Это два деления с остатком в $\Om[X]$ (во втором остаток равен $0$), и по единственности (\rf{1.6}) $r=0$.
\end{proof}

\subsection*{Определение кратности и её независимость от поля}

Пусть $f\in F[X]$, $f\ne0$, и $\alpha$~--- корень $f$ в некотором расширении $E\supset F$. Наибольшее $m\ge1$, для
которого $(X-\alpha)^m$ делит $f$ в $E[X]$, называется \emph{кратностью} корня $\alpha$ (такое $m$ существует, так как
$m\le\deg f$). Иначе говоря, $m$~--- это показатель в записи
\[
f(X)=(X-\alpha)^mg(X),\qquad g\in E[X],\quad g(\alpha)\ne0
\]
(если бы $g(\alpha)=0$, то $X-\alpha$ делил бы $g$ по \rf{1.7}, и $(X-\alpha)^{m+1}$ делил бы $f$). Кратность не зависит от
того, в каком поле, содержащем $\alpha$, мы её вычисляем: если $\alpha\in E\subset E'$, то по \rf{2.17a} (применённой к
$E\subset E'$ и многочленам $(X-\alpha)^k\in E[X]$) делимость $f$ на $(X-\alpha)^k$ в $E[X]$ и в $E'[X]$~--- одно и то же.
Корень кратности $1$ называется \emph{простым}, корень кратности $m>1$~--- \emph{кратным}.

По \rf{2.7} у $f$ есть поле разложения $E$, и в $E[X]$
\begin{equation}\label{eq:ch2B-split}
f(X)=a\prod_{i=1}^r(X-\alpha_i)^{m_i},\qquad a\in F,\quad \alpha_i\ \text{попарно различны},\quad m_i\ge1,\quad
\sum_{i=1}^rm_i=\deg f .
\end{equation}
Здесь $m_i$~--- в точности кратность корня $\alpha_i$: $(X-\alpha_i)^{m_i}$ делит $f$, а частное
$a\prod_{j\ne i}(X-\alpha_j)^{m_j}$ в точке $\alpha_i$ не обращается в нуль.

Неупорядоченный набор кратностей $m_1,\dots,m_r$ в \eqref{eq:ch2B-split} \emph{не зависит от выбора поля}, в котором $f$
раскладывается на линейные множители. Если $f$ раскладывается в расширении $E$, то все $\alpha_i$ и $a$ лежат в подполе
$F[\alpha_1,\dots,\alpha_r]$, равенство \eqref{eq:ch2B-split} выполняется уже в $F[\alpha_1,\dots,\alpha_r][X]$, а
кратности, как мы видели, от объемлющего поля не зависят; поэтому можно считать, что $E$~--- поле разложения $f$. Пусть
$E$ и $E'$~--- два поля разложения $f$, и
\[
f(X)=a\prod_{i=1}^r(X-\alpha_i)^{m_i}\ \text{в}\ E[X],\qquad f(X)=a\prod_{j=1}^{r'}(X-\alpha'_j)^{m'_j}\ \text{в}\ E'[X].
\]
По \rf{2.13} существует $F$-изоморфизм $\varphi\colon E\to E'$; продолжим его до изоморфизма колец $E[X]\to E'[X]$,
применяя $\varphi$ к коэффициентам ($X\mapsto X$). Он тождествен на $F[X]$, поэтому переводит $f$ в $f$, а разложение $f$
в $E[X]$~--- в разложение
\[
f(X)=a\prod_{i=1}^r\bigl(X-\varphi(\alpha_i)\bigr)^{m_i}\quad\text{в}\ E'[X],
\]
где $\varphi(\alpha_i)$ попарно различны. Кольцо $E'[X]$ факториально, а многочлены $X-\beta$ в нём неприводимы и при
разных $\beta$ не ассоциированы; поэтому это разложение совпадает с разложением через $\alpha'_j$ с точностью до
нумерации. Следовательно, $r=r'$ и
\[
\{m_1,\dots,m_r\}=\{m'_1,\dots,m'_{r'}\}\qquad(\text{равенство мультимножеств}).
\]

\par\noindent\begin{minipage}{\linewidth}
\[
\begin{tikzcd}[column sep=large,row sep=normal]
& E\arrow[dd,cnew,"\varphi","\sim"'] & E[X]\arrow[dd,cnew,"{\varphi,\ X\,\mapsto\,X}"] & a\prod(X-\alpha_i)^{m_i}\arrow[dd,cnew,mapsto]\\
F\arrow[ur,cold,hook]\arrow[dr,cold,hook] & & &\\
& E' & E'[X] & a\prod\bigl(X-\varphi(\alpha_i)\bigr)^{m_i}=a\prod(X-\alpha'_j)^{m'_j}
\end{tikzcd}
\]
\figcap{$F$-изоморфизм полей разложения $\varphi$ переносит разложение $f$ на линейные множители в $E[X]$ в разложение
в $E'[X]$; по однозначности разложения кратности совпадают.}
\end{minipage}\par\medskip

Говорят, что $f$ \emph{имеет кратный корень}, если хотя бы одно $m_i>1$, и что $f$ \emph{имеет только простые корни},
если все $m_i=1$. Итак, <<$f$ имеет кратный корень>> означает <<$f$ имеет кратный корень в одном, а значит, и в каждом
расширении $F$, в котором $f$ раскладывается на линейные множители>>.

\section{Сепарабельные многочлены}

Когда многочлен имеет кратный корень? Если $f$ имеет кратный множитель в $F[X]$, скажем $f=g^2h$ с непостоянным $g$, то,
очевидно, у $f$ есть кратный корень (любой корень $g$). Если же $f$~--- произведение \emph{различных} унитарных
неприводимых многочленов (и константы), то по \rf{2.17} корни разных множителей различны, и $f$ имеет кратный корень
тогда и только тогда, когда кратный корень есть у одного из множителей. Значит, достаточно понять, когда кратный корень
есть у \emph{неприводимого} многочлена.

\begin{example}[2.18]
Пусть $\operatorname{char}F=p\ne0$, и $F$ содержит элемент $a$, не являющийся $p$-й степенью, например $a=T$ в поле
$\F_p(T)$ (\rf{1.4a}). Тогда $X^p-a$ неприводим в $F[X]$ (\rf{2.3}), но в его поле разложения $X^p-a=(X-\alpha)^p$, где
$\alpha^p=a$ (Фробениус аддитивен, см.\ раздел <<Эндоморфизм Фробениуса>> перед \rf{1.4a}). Итак, неприводимый многочлен
может иметь кратные корни: у $X^p-T$ над $\F_p(T)$ единственный корень кратности $p$.
\end{example}

Производная $f'$ многочлена $f=\sum a_iX^i$ определена в разделе <<Производные и формула Тейлора>> перед \rf{1.7a}:
$f'=\sum ia_iX^{i-1}$. Она не зависит от поля, над которым рассматривается $f$ (коэффициенты те же). Обычные правила
дифференцирования сохраняются.

\begin{lemma}[2.18a]
Для $f,g\in F[X]$ и $c\in F$: $(f+g)'=f'+g'$, $(cf)'=cf'$, $(fg)'=f'g+fg'$ и $(g^m)'=mg^{m-1}g'$ при $m\ge1$.
\end{lemma}

\begin{proof}
Линейность видна из определения. Обе части правила Лейбница билинейны по $(f,g)$, поэтому его достаточно проверить для
одночленов: $(X^{i+j})'=(i+j)X^{i+j-1}=iX^{i-1}\cdot X^j+X^i\cdot jX^{j-1}$. Последнее правило~--- индукция по $m$:
$(g^{m+1})'=(g^m)'g+g^mg'=mg^{m-1}g'g+g^mg'=(m+1)g^mg'$.
\end{proof}

Подчеркнём, что в характеристике $p$ производная $X^p$ равна $pX^{p-1}=0$: ненулевой многочлен может иметь нулевую
производную.

\begin{lemma}[2.19]
Корень $\alpha$ многочлена $f\ne0$ является кратным тогда и только тогда, когда он является также корнем $f'$.
\end{lemma}

\begin{proof}
Пусть $\alpha$ лежит в расширении $E\supset F$ (любом, содержащем $\alpha$), и
\[
f(X)=(X-\alpha)^mg(X),\qquad m\ge1,\quad g\in E[X],\quad g(\alpha)\ne0,
\]
где $m$~--- кратность $\alpha$. По \rf{2.18a}
\begin{equation}\label{eq:ch2B-deriv}
f'(X)=\begin{cases}g(X)+(X-\alpha)g'(X),& m=1,\\[2pt]
m(X-\alpha)^{m-1}g(X)+(X-\alpha)^mg'(X),& m>1.\end{cases}
\end{equation}
Поэтому $f'(\alpha)=g(\alpha)\ne0$ при $m=1$ и $f'(\alpha)=0$ при $m>1$.
\end{proof}

\begin{proposition}[2.20]
Для непостоянного неприводимого многочлена $f\in F[X]$ следующие условия равносильны:
\begin{enumerate}[label=(\alph*)]
\item $f$ имеет кратный корень;
\item $\gcd(f,f')\ne1$;
\item $F$ имеет ненулевую характеристику $p$, и $f$~--- многочлен от $X^p$;
\item все корни $f$ кратные.
\end{enumerate}
\end{proposition}

\begin{proof}
(a)$\Rightarrow$(b). Если $\alpha$~--- кратный корень $f$ (в каком-либо расширении $E$), то по \rf{2.19} $X-\alpha$~---
общий делитель $f$ и $f'$ в $E[X]$. Значит, $\gcd(f,f')\ne1$ в $E[X]$, а тогда по \rf{2.17} и в $F[X]$.

(b)$\Rightarrow$(c). Многочлен $d=\gcd(f,f')$ делит неприводимый $f$ и не является константой, поэтому $d=cf$,
$c\in F^\times$. Значит, $f$ делит $f'$; так как $\deg f'<\deg f$, получаем $f'=0$. Пусть $f=a_0+a_1X+\dots+a_dX^d$,
$d\ge1$, $a_d\ne0$. Тогда
\[
f'=a_1+2a_2X+\dots+ia_iX^{i-1}+\dots+da_dX^{d-1},
\]
и $f'=0$ означает $ia_i=0$ для всех $i\ge1$. В характеристике $0$ это дало бы $a_i=0$ при всех $i\ge1$, в частности
$a_d=0$,~--- противоречие. Значит, $\operatorname{char}F=p\ne0$, и $ia_i=0$ означает $a_i=0$ для всех $i$, не делящихся на
$p$ (при $p\nmid i$ элемент $i\cdot1_F$ обратим). Итак, $f=a_0+a_pX^p+a_{2p}X^{2p}+\dots$~--- многочлен от $X^p$.

(c)$\Rightarrow$(d). По условию $f(X)=g(X^p)$, $g\in F[X]$. Пусть $g(X)=c\prod_{i=1}^r(X-a_i)^{m_i}$ с попарно
различными $a_i$ в некотором расширении $E_1$ поля $F$ (\rf{2.7}), и пусть $E_2\supset E_1$~--- расширение, в котором
раскладывается $\prod_i(X^p-a_i)$; в нём $a_i=\alpha_i^p$ для некоторых $\alpha_i$. Тогда
$X^p-a_i=X^p-\alpha_i^p=(X-\alpha_i)^p$ и
\[
f(X)=g(X^p)=c\prod_{i=1}^r(X^p-a_i)^{m_i}=c\prod_{i=1}^r(X-\alpha_i)^{pm_i},
\]
причём $\alpha_i$ попарно различны, так как различны их $p$-е степени $a_i$. Значит, каждый корень $f$ имеет кратность
$pm_i\ge p\ge2$.

(d)$\Rightarrow$(a). Очевидно: у непостоянного $f$ корни есть (\rf{2.7}).
\end{proof}

Заметим, что для корня $\alpha$ кратности $m$ многочлен $(X-\alpha)^{m-1}$ делит $\gcd(f,f')$ (по
\eqref{eq:ch2B-deriv}), но, вообще говоря, не равен ему: для $f=(X-1)^2(X-2)^2\in\Q[X]$ имеем $f'=2(X-1)(X-2)(2X-3)$ и
$\gcd(f,f')=(X-1)(X-2)$, а не $X-1$. Для неприводимого $f$ с кратным корнем, как показывает (b)$\Rightarrow$(c), всегда
$f'=0$ и $\gcd(f,f')=c^{-1}f$, где $c$~--- старший коэффициент $f$.

\begin{example}[2.20a]
Пусть $\operatorname{char}k=p$, $F=k(t)$, и $f=X^{2p}+tX^p+t$; этот многочлен неприводим над $F$ по критерию Эйзенштейна
(\rf{1.16d}(b)). Он является многочленом от $X^p$: $f=g(X^p)$, $g=X^2+tX+t$. Многочлен $g$ имеет различные корни $b_1\ne
b_2$: при $p=2$ его производная равна ненулевой константе $t$, а при $p\ne2$ его дискриминант $t^2-4t=t(t-4)$ не равен
нулю. Как в доказательстве (c)$\Rightarrow$(d), $f=(X-\beta_1)^p(X-\beta_2)^p$, где $\beta_i^p=b_i$: у $f$ ровно два
различных корня, каждый кратности $p$. Общее утверждение такого рода~--- упражнение 2-6.
\end{example}

\begin{proposition}[2.21]
Для ненулевого многочлена $f\in F[X]$ следующие условия равносильны:
\begin{enumerate}[label=(\alph*)]
\item $\gcd(f,f')=1$ в $F[X]$;
\item $f$ имеет только простые корни.
\end{enumerate}
\end{proposition}

\begin{proof}
Пусть $\Om$~--- расширение $F$, в котором $f$ раскладывается на линейные множители (\rf{2.7}). По \rf{2.19} корень $\alpha\in\Om$
многочлена $f$ кратен тогда и только тогда, когда он является и корнем $f'$.

(a)$\Rightarrow$(b). Если $\gcd(f,f')=1$, то по \rf{2.17} $f$ и $f'$ не имеют общих множителей в $\Om[X]$, в частности
общих корней; значит, все корни $f$ простые.

(b)$\Rightarrow$(a). Если $d=\gcd(f,f')$ непостоянен, то он делит $f$ и потому раскладывается в $\Om$ на линейные
множители; любой его корень $\alpha\in\Om$~--- общий корень $f$ и $f'$, т.\,е.\ кратный корень $f$.
\end{proof}

Условие (a) проверяется алгоритмом Евклида в $F[X]$, без нахождения корней. Например, $f=X^3-3X+2$ над $\Q$:
$f'=3(X-1)(X+1)$, $\gcd(f,f')=X-1$, и действительно $f=(X-1)^2(X+2)$. Над $\F_3$ тот же многочлен равен $X^3-1=(X-1)^3$, а
его производная равна нулю.

\begin{definition}[2.22]
Многочлен называется \emph{сепарабельным}, если он ненулевой и удовлетворяет равносильным условиям \rf{2.21}.
\end{definition}

Перечислим простые следствия.
\begin{enumerate}[label=(\alph*)]
\item Ненулевые константы сепарабельны ($\gcd(c,0)=1$; корней нет).
\item Непостоянный неприводимый $f$ сепарабелен, кроме случая, когда $\operatorname{char}F=p\ne0$ и $f$~--- многочлен от
$X^p$ (\rf{2.20}). В частности, в характеристике $0$ все неприводимые многочлены сепарабельны, а в характеристике $p$
неприводимый $f$ сепарабелен, если $p\nmid\deg f$ (многочлен от $X^p$ имеет степень, кратную $p$).
\item Пусть $f=c\prod_if_i$, где $c\in F^\times$, а $f_i$~--- унитарные неприводимые (не обязательно различные). Тогда $f$
сепарабелен $\iff$ все $f_i$ попарно различны и сепарабельны. Действительно, если $f_i=f_j$ при $i\ne j$, то любой корень
$f_i$~--- кратный корень $f$; если $f_i$ имеет кратный корень, то это и кратный корень $f$. Обратно, если $f_i$ попарно
различны и сепарабельны, то корни каждого $f_i$ простые, а корни разных $f_i$ различны (\rf{2.17}); значит, в поле, где
$f$ раскладывается, все линейные множители $f$ различны.
\item Если $f$ сепарабелен как многочлен из $F[X]$, то он сепарабелен и как многочлен из $E[X]$ для любого расширения $E$
поля $F$: $f'$ один и тот же, а $\gcd(f,f')=1$ по \rf{2.17}. (Неприводимые множители $f$ над $E$ могут быть другими,
но они различны и сепарабельны по (c).)
\end{enumerate}

\begin{remark}[2.22a]
Иногда (особенно в старой литературе) сепарабельным называют многочлен, каждый \emph{неприводимый множитель} которого
имеет только простые корни. При таком определении, например, $(X-1)^2\in\Q[X]$ сепарабелен, а в нашем смысле~--- нет.
Для неприводимых многочленов определения совпадают, а именно они нужны в определении сепарабельного расширения (ниже и в
главе 3). Наше определение удобнее тем, что сепарабельность многочлена~--- свойство его корней и не зависит от поля
(пункт (d)). При старом определении это не так: $X^p-T$ над $\F_p(T)$ несепарабелен (его единственный неприводимый
множитель $X^p-T$ имеет кратный корень), а над $E=\F_p(T)[T^{1/p}]$ он равен $(X-T^{1/p})^p$, и его неприводимый
множитель $X-T^{1/p}$ имеет простой корень.
\end{remark}

\subsection*{Неприводимость многочлена Артина--Шрайера}

Многочлен $X^p-X-a$ в характеристике $p$ имеет степень $p$, но сепарабелен всегда: его производная равна $pX^{p-1}-1=-1$,
и $\gcd(f,-1)=1$. Его неприводимость проверяется особенно просто: такой многочлен либо неприводим, либо раскладывается
полностью.

\begin{proposition}[2.22b: многочлен Артина--Шрайера]
Пусть $\operatorname{char}F=p\ne0$, $a\in F$, $f=X^p-X-a$.
\begin{enumerate}[label=(\alph*)]
\item Если $\alpha$~--- корень $f$ в каком-либо расширении $F$, то корни $f$~--- это $p$ различных элементов $\alpha+i$,
$i\in\F_p$, и $F[\alpha]$~--- поле разложения $f$.
\item Если $a=\beta^p-\beta$ для некоторого $\beta\in F$, то $f$ раскладывается в $F[X]$ в произведение $p$ различных
линейных множителей $X-\beta-i$, $i\in\F_p$. В противном случае $f$ неприводим в $F[X]$.
\end{enumerate}
\end{proposition}

\begin{proof}
(a) Это пример \rf{2.9}: $f=\prod_{i\in\F_p}(X-\alpha-i)$.

(b) Первое утверждение~--- это (a) при $\alpha=\beta$. Пусть $f$ приводим: $f=gh$, где $g\in F[X]$ унитарен и
$0<r=\deg g<p$. Пусть $\alpha$~--- корень $f$ в некотором расширении $F$ (\rf{2.7}). По (a) и однозначности разложения в $F[\alpha][X]$
\[
g=\prod_{i\in I}(X-\alpha-i)\qquad\text{для некоторого }I\subset\F_p,\ |I|=r .
\]
Коэффициент $g$ при $X^{r-1}$ равен $-\sum_{i\in I}(\alpha+i)=-r\alpha-\sum_{i\in I}i$ и лежит в $F$. Так как $0<r<p$,
элемент $r\cdot1_F$ обратим, и $\alpha\in F$. Значит, $a=\alpha^p-\alpha$ с $\alpha\in F$.
\end{proof}

Итак, $X^p-X-a$ неприводим над $F$ тогда и только тогда, когда $a\notin\{\beta^p-\beta\mid\beta\in F\}$; а если он
приводим, то сразу раскладывается на различные линейные множители. Например, над $\F_p$ многочлен $X^p-X-a$ при $a\ne0$
неприводим: для $\beta\in\F_p$ имеем $\beta^p-\beta=0$. Расширение $F[\alpha]/F$, порождённое корнем неприводимого
многочлена $X^p-X-a$, циклично степени $p$ с образующей $\alpha\mapsto\alpha+1$, и все циклические расширения степени
$p$ в характеристике $p$ так получаются; это теория Артина--Шрайера (\rf{5.32c}); критерий неприводимости над конечными
полями через след~--- в \rf{5.32f}, примеры над $k(t)$~--- в \rf{5.32h}.

\begin{aside}[2.22c: формулы Виета для $X^p-X$]
Все $p$ элементов $\F_p$~--- корни $X^p-X$ (малая теорема Ферма), поэтому
\[
X^p-X=\prod_{k\in\F_p}(X-k)\qquad\text{в }\F_p[X].
\]
Сравнивая коэффициенты, получаем: элементарные симметрические многочлены от всех элементов $\F_p$ степеней
$1,\dots,p-2$ равны нулю (например, $0+1+\dots+(p-1)=0$ при $p>2$), а коэффициент при $X$ даёт
$(-1)^{p-1}\prod_{k\ne0}k=-1$ (единственное произведение $p-1$ элементов без нуля~--- произведение всех ненулевых), т.\,е.
\[
\prod_{k\in\F_p^\times}k=-1,\qquad\text{или}\qquad (p-1)!\equiv-1\pmod p
\]
(теорема Вильсона; при $p=2$ обе части равны $1=-1$). Подставив $X-\alpha$ вместо $X$, получаем тождество
\[
\prod_{k\in\F_p}(X-\alpha-k)=(X-\alpha)^p-(X-\alpha)=X^p-X-(\alpha^p-\alpha),
\]
т.\,е.\ ещё раз получаем разложение из \rf{2.9}. Формулы Виета для корней $\alpha+k$ многочлена $X^p-X-a$ дают
$\sum_k(\alpha+k)=0$ при $p>2$ и $\prod_k(\alpha+k)=a$ (при подстановке $X=0$ получаем
$(-1)^p\prod_k(\alpha+k)=-a$, а $(-1)^{p+1}=1$ при нечётном $p$, и $-a=a$ при $p=2$). Вывести отсюда, что $a$ имеет вид
$\beta^p-\beta$ с $\beta\in F$, нельзя (эти равенства верны для любого корня $\alpha$, лежащего в $F$ или нет); для этого
и нужен аргумент с коэффициентом делителя $g$ в \rf{2.22b}(b).
\end{aside}

\subsection*{Сепарабельные элементы и простые расширения}

\begin{definition}[2.22d]
Пусть $E/F$~--- алгебраическое расширение. Элемент $\alpha\in E$ \emph{сепарабелен над $F$}, если его минимальный
многочлен над $F$ сепарабелен. Расширение $E/F$ \emph{сепарабельно}, если каждый элемент $E$ сепарабелен над $F$
(подробнее~--- в главе 3, \rf{3.6}, \rf{3.15}).
\end{definition}

В характеристике $0$ каждое алгебраическое расширение сепарабельно (по пункту (b) после \rf{2.22}). Естественный вопрос:
если $\alpha$ сепарабелен над $F$, то сепарабельно ли всё расширение $F[\alpha]/F$, т.\,е.\ сепарабелен ли каждый
многочлен от $\alpha$? Ответ положителен; в главе 3 это будет получено из теории Галуа (\rf{3.15}), а здесь мы дадим
прямое доказательство. (Обратное тривиально: если $F[\alpha]/F$ сепарабельно, то сепарабелен и элемент $\alpha$.)

Пусть $\operatorname{char}F=p\ne0$ и $E/F$ конечно. Положим $E^p=\{x^p\mid x\in E\}$; это подполе $E$ (образ
Фробениуса), и через $F[E^p]$ обозначим подкольцо $E$, порождённое $F$ и $E^p$. Произведение $p$-х степеней~---
$p$-я степень, поэтому (\rf{1.21}) $F[E^p]$ состоит из линейных комбинаций элементов $E^p$ с коэффициентами из $F$, т.\,е.\
$F[E^p]$~--- это $F$-линейная оболочка $E^p$; по \rf{1.23} это подполе $E$.

\begin{lemma}[2.22e]
Пусть $\operatorname{char}F=p\ne0$, и $\beta$ (алгебраичный над $F$) сепарабелен над $F$. Тогда $\beta\in F[\beta^p]$.
\end{lemma}

\begin{proof}
Положим $K=F[\beta^p]$, и пусть $h$~--- минимальный многочлен $\beta$ над $K$. Так как $\beta^p\in K$, то $\beta$~---
корень многочлена $X^p-\beta^p=(X-\beta)^p\in K[X]$, и $h$ делит $(X-\beta)^p$; по однозначности разложения в
$K[\beta][X]$ имеем $h=(X-\beta)^s$, $s\ge1$. С другой стороны, $h$ делит минимальный многочлен $f$ элемента $\beta$ над
$F$ (он лежит в $K[X]$ и имеет корень $\beta$), а у $f$ все корни простые. Значит, $s=1$, $h=X-\beta\in K[X]$, т.\,е.\
$\beta\in K$.
\end{proof}

\begin{theorem}[2.22f]
Пусть $\operatorname{char}F=p\ne0$, и $E/F$ конечно. Расширение $E/F$ сепарабельно тогда и только тогда, когда
$E=F[E^p]$.
\end{theorem}

\begin{proof}
$\Rightarrow$. Каждый $\beta\in E$ сепарабелен, и по \rf{2.22e} $\beta\in F[\beta^p]\subset F[E^p]$.

$\Leftarrow$. Сначала покажем: \emph{если $x_1,\dots,x_m\in E$ линейно независимы над $F$, то и $x_1^p,\dots,x_m^p$
линейно независимы над $F$.} Дополним $x_1,\dots,x_m$ до базиса $x_1,\dots,x_n$ пространства $E$ над $F$. Любой элемент
$E^p$ имеет вид $\bigl(\sum c_ix_i\bigr)^p=\sum c_i^px_i^p$, $c_i\in F$, т.\,е.\ лежит в $F$-линейной оболочке $V$
элементов $x_1^p,\dots,x_n^p$. Значит, $E=F[E^p]$ (линейная оболочка $E^p$) содержится в $V$, и $n$ элементов
$x_1^p,\dots,x_n^p$ порождают $n$-мерное пространство $E$; поэтому они образуют базис и, в частности, линейно
независимы.

Пусть теперь $\beta\in E$, и $f$~--- его минимальный многочлен над $F$, $d=\deg f$. Элементы $1,\beta,\dots,\beta^{d-1}$
линейно независимы над $F$, а потому независимы и $1,\beta^p,\dots,\beta^{p(d-1)}$, т.\,е.\ степень $\beta^p$ над $F$ не
меньше $d$. Если бы $f$ был несепарабелен, то $f=g(X^p)$ (\rf{2.20}), $g(\beta^p)=0$ и $\deg g=d/p<d$~--- противоречие.
Значит, $\beta$ сепарабелен над $F$.
\end{proof}

\begin{corollary}[2.22g]
Если $E=F[\alpha_1,\dots,\alpha_n]$, где элементы $\alpha_i$ алгебраичны и сепарабельны над $F$, то расширение $E/F$
сепарабельно. В частности, простое расширение $F[\alpha]$, порождённое сепарабельным элементом $\alpha$, сепарабельно.
\end{corollary}

\begin{proof}
В характеристике $0$ доказывать нечего. Пусть $\operatorname{char}F=p$. Расширение $E/F$ конечно (\rf{1.30}), и по \rf{2.22e}
$\alpha_i\in F[\alpha_i^p]\subset F[E^p]$ для всех $i$. Поэтому $E=F[\alpha_1,\dots,\alpha_n]\subset F[E^p]\subset E$, и
по \rf{2.22f} $E/F$ сепарабельно.
\end{proof}

Например, $F=\F_p(T)$ и $\alpha$~--- корень $X^p-X-T$. Этот многочлен неприводим по \rf{2.22b}, так как $T$ не имеет
вида $\beta^p-\beta$ с $\beta\in\F_p(T)$. Действительно, запишем $\beta=g/h$, где $g,h\in\F_p[T]$ взаимно просты и $h$
унитарен. Если $h$ непостоянен, то $\beta^p-\beta=(g^p-gh^{p-1})/h^p$, и числитель взаимно прост с $h$ (по модулю любого
неприводимого делителя $h$ он сравним с $g^p$), так что $\beta^p-\beta$~--- не многочлен; если $\beta$~--- многочлен
степени $n\ge1$, то $\deg(\beta^p-\beta)=pn\ne1$; если $\beta\in\F_p$, то $\beta^p-\beta=0$ (ср.\ \rf{5.32h}). Кроме того,
$X^p-X-T$ сепарабелен (его производная равна $-1$), так что каждый элемент поля $F[\alpha]$ сепарабелен над $F$.
Напротив, $F[\alpha]$ с $\alpha^p=T$ (пример \rf{2.18}) не сепарабельно уже потому, что несепарабелен сам $\alpha$.

\begin{remark}[2.22h: транзитивность]
Критерий \rf{2.22f} сразу даёт транзитивность сепарабельности для конечных расширений. Пусть $F\subset E\subset K$,
$K/E$ и $E/F$ конечны и сепарабельны, $\operatorname{char}F=p$. Тогда $K=E[K^p]$ и $E=F[E^p]$, откуда
\[
K=E[K^p]=F[E^p][K^p]=F[K^p],
\]
так как $E^p\subset K^p$; значит, $K/F$ сепарабельно. Для произвольных алгебраических расширений это утверждение (оно
сводится к конечному случаю) доказано в \rf{6.8b}.
\end{remark}

\section{Совершенные поля}

\begin{definition}[2.23]
Поле $F$ называется \emph{совершенным}, если оно имеет характеристику $0$ либо имеет характеристику $p$ и каждый
элемент $F$ является $p$-й степенью (т.\,е.\ из каждого элемента $F$ можно извлечь в $F$ корень степени $p$).
\end{definition}

Иначе говоря, поле $F$ с характеристической экспонентой $q$ (см.\ после \rf{1.4a}) совершенно, если $F=F^q$, т.\,е.\
если эндоморфизм $x\mapsto x^q$ поля $F$ сюръективен (а значит, является автоморфизмом).

\begin{proposition}[2.24]
Поле $F$ совершенно тогда и только тогда, когда каждый неприводимый многочлен из $F[X]$ сепарабелен.
\end{proposition}

\begin{proof}
В характеристике $0$ оба условия выполнены (пункт (b) после \rf{2.22}). Пусть $\operatorname{char}F=p\ne0$. Если $F$
содержит элемент $a$, не являющийся $p$-й степенью, то $X^p-a$ неприводим в $F[X]$, но не сепарабелен (\rf{2.18}).
Обратно, пусть каждый элемент $F$ является $p$-й степенью, и $f\in F[X]$ неприводим. Если бы $f$ был несепарабелен, то
по \rf{2.20} $f=\sum_ia_iX^{ip}$; запишем $a_i=b_i^p$, $b_i\in F$. Так как возведение в степень $p$ в $F[X]$~---
гомоморфизм колец,
\[
\sum_ia_iX^{ip}=\sum_ib_i^p\bigl(X^i\bigr)^p=\Bigl(\sum_ib_iX^i\Bigr)^p,
\]
т.\,е.\ $f$~--- $p$-я степень непостоянного многочлена и не может быть неприводимым. Значит, $f$ сепарабелен.
\end{proof}

\begin{example}[2.25]
\begin{enumerate}[label=(\alph*)]
\item Конечное поле $F$ совершенно: эндоморфизм Фробениуса $a\mapsto a^p\colon F\to F$ инъективен (как любой
гомоморфизм полей: $a^p=0\Rightarrow a=0$), а значит, и сюръективен, так как $F$ конечно.
\item Поле, являющееся объединением возрастающей цепочки (или, более общо, направленного семейства) совершенных
подполей, совершенно: каждый элемент лежит в одном из этих подполей и является в нём $p$-й степенью. Каждое поле $E$,
алгебраическое над $\F_p$, совершенно. Здесь удобнее рассуждать поэлементно, не выписывая цепочку подполей (она
может оказаться несчётной): для $a\in E$ поле $\F_p[a]$ конечно над $\F_p$ (\rf{1.30}), т.\,е.\ является конечным полем, и
по (a) $a=b^p$ для некоторого $b\in\F_p[a]\subset E$.
\item Каждое алгебраически замкнутое поле $F$ совершенно: у многочлена $X^p-a$ есть корень в $F$.
\item Если $\operatorname{char}F_0=p\ne0$, то поле $F=F_0(X)$ рациональных функций несовершенно: $X$ не является $p$-й
степенью. Действительно, из $X=(g/h)^p$ с $g,h\in F_0[X]$, $h\ne0$, следовало бы $Xh^p=g^p$ и
$1+p\deg h=p\deg g$, что невозможно (ср.\ \rf{1.4a}).
\end{enumerate}
\end{example}

\begin{proposition}[2.25a]
Каждое алгебраическое расширение $E$ совершенного поля $F$ совершенно.
\end{proposition}

\begin{proof}
В характеристике $0$ доказывать нечего; пусть $\operatorname{char}F=p$ и $a\in E$. Поле $E_0=F[a]$ конечно над
$F$ и сепарабельно над ним (\rf{2.24}: минимальные многочлены над $F$ неприводимы, а значит, сепарабельны). По
\rf{2.22f} $E_0=F[E_0^p]$, а так как $F=F^p\subset E_0^p$, то $F[E_0^p]=E_0^p$. Итак, $E_0=E_0^p$, и $a$~--- $p$-я степень
элемента $E_0\subset E$.
\end{proof}

\begin{example}[2.25b]
Совершенное поле может быть объединением возрастающей цепочки \emph{несовершенных} полей. Пусть $\Om$~--- алгебраическое
замыкание поля $\F_p(T)$, $T_n\in\Om$~--- корень $X^{p^n}-T$ (он единствен, так как $X^{p^n}-T=(X-T_n)^{p^n}$, и
$T_{n+1}^p=T_n$), и $K_n=\F_p(T_n)$. Ни одно $K_n$ не совершенно (\rf{2.25}(d): $T_n$ трансцендентен над $\F_p$), но
$K_0\subset K_1\subset\dots$ и Фробениус переводит $K_{n+1}=\F_p(T_{n+1})$ на $\F_p(T_{n+1}^p)=K_n$. Поэтому каждый
элемент $K_n$ является $p$-й степенью элемента $K_{n+1}$, и объединение $\bigcup_nK_n=\F_p(T^{1/p^\infty})$ совершенно.
\end{example}

\begin{aside}[2.26]
Когда $F$ совершенно, мы увидим (\rf{5.1}), что каждое конечное расширение $E/F$ простое, т.\,е.\ $E=F[\alpha]$, где
$\alpha$~--- корень (сепарабельного) многочлена $f\in F[X]$ степени $[E:F]$. Тогда из \rf{2.1}(b) сразу следует, что для
любого расширения $\Om$ поля $F$ число $F$-гомоморфизмов $E\to\Om$ не превосходит $[E:F]$, причём равенство имеет место
тогда и только тогда, когда $f$ раскладывается в $\Om$ на линейные множители. Здесь этим рассуждением воспользоваться
нельзя: теорема о примитивном элементе \rf{5.1} доказывается с помощью результатов этой и следующей глав, и мы получили
бы порочный круг.
\end{aside}
